
\subsection{Vetor}

É uma quantidade com direção, sentido e magnitude (tamanho).

\emph{Exemplo:}
a velocidade de uma espaçonave não é apenas $20 [\frac{m}{s}]$, mas $20 [\frac{m}{s}]$ para o norte.

\begin{itemize}
    \item Magnitude: módulo $(20 [\frac{m}{s}])$
    \item Direção: linha para qual aponta (norte-sul)
    \item Sentido: seta para o norte (e não para o sul)
\end{itemize}

\subsubsection{Representação matemática}
Pode ser escrito como uma seta ou como uma coluna de números (coordenadas) em um sistema cartesiano.

\begin{equation}
\label{eq:representacao_vetor1}
    \vec v =
    \begin{bmatrix}
    v_x \\ v_y \\ v_z
    \end{bmatrix}
\end{equation}

Por exemplo:
\begin{equation}
\label{eq:exemplo_vetor1}
    \vec v =
    \begin{bmatrix}
    3 \\ -2 \\ 5
    \end{bmatrix}
\end{equation}

\subsubsection{Operações básicas com vetores: Adição}
Soma componente a componente.

\begin{equation}
\label{eq:soma_vetores}
    \vec u + \vec v =
    \begin{bmatrix}
    u_x+v_x \\
    u_y+v_y \\
    u_z+v_z
    \end{bmatrix}
\end{equation}

\subsubsection{Operações básicas com vetores: Multiplicação escalar}
Multiplica um número escalar (sem direção ou sentido, por exemplo) pela matriz em questão.

\begin{equation}
\label{eq:multiplicacao_escalar}
    \vec k * \vec v = 
    \begin{bmatrix}
    k\vec v_x \\
    k\vec v_y \\
    k\vec v_z
    \end{bmatrix}
\end{equation}

\subsubsection{Operações básicas com vetores: Transposta}
Altera a posição dos termos, caso seja matriz $linha$, ela se torna matriz $coluna$ (vetor).

\begin{equation}
\label{eq:transposta_linha2coluna1}
    \vec x = 
    [x,y,z]
\end{equation}

\begin{equation}
\label{eq:transposta_linha2coluna2}
    \vec x^T =
    \begin{bmatrix}
    x \\
    y \\
    z
    \end{bmatrix} 
\end{equation}

\subsubsection{Tipos de matrizes: Quadradas}
Podem ser $2\times2, 3\times3, 4\times4, ..., n\times n$, ou seja, elas possuem o mesmo número de colunas e linhas. Normalmente são chamadas de matriz de ordem $n$; onde $n$ representa o numero de linhas (ou colunas). Uma matriz diagonal \textbf{A} de ordem $n$ é normalmente descrita como:


\subsection{Aplicação em robótica ou espaço}
O vetor pode representar uma posição no espaço (ex.: posição de um satélite) ou também pode representar velocidade, força ou até um eixo de rotação.

